\section{引言：为什么要讲 policy gradient}

这一讲是本课程的最后一讲，由于前一讲已经概览了 \emph{reinforcement learning from verifiable rewards} 和策略梯度类方法，这一讲不再引入全新框架，而是把 policy gradient 和 GRPO 的机制掰开揉碎，重点看三个问题：
\begin{enumerate}
    \item 在语言模型里，RL 的 \textbf{state}、\textbf{action}、\textbf{reward} 到底对应什么；
    \item 为什么 naive policy gradient 会有高方差，baseline 和 advantage 为什么能缓解它；
    \item GRPO 在代码里到底怎么跑，为什么要用 group 内平均奖励做比较。
\end{enumerate}

\begin{importantbox}{这一讲的核心}
\begin{itemize}
    \item \textbf{如果你能度量，就能优化}，但 RL 的难点在于度量很稀疏、很噪声。
    \item \textbf{策略梯度本身不神秘}，本质上就是对 $\log \pi(a \mid s)$ 乘上一个奖励信号。
    \item \textbf{baseline 的作用是降方差}，不是改变优化目标。
\end{itemize}
\end{importantbox}


\begin{knowledgebox}{为什么这一讲很像“工程课”}
因为 policy gradient 的关键问题不是能不能写公式，而是如何把公式变成一个在 `verifiable reward` 下能稳定训练的系统。只要 reward 稍微稀疏一点，方差、采样、baseline、clipping、reference model 都会立刻变成工程问题。
\end{knowledgebox}

